\documentclass[11pt]{amsart}
\usepackage[margin=1in]{geometry}
\usepackage{amsmath,amssymb,amsthm,mathtools}
\usepackage{booktabs}
\usepackage{enumitem}
\usepackage[colorlinks=true,linkcolor=blue,citecolor=blue,urlcolor=blue]{hyperref}

\newtheorem{theorem}{Theorem}[section]
\newtheorem{lemma}[theorem]{Lemma}
\newtheorem{proposition}[theorem]{Proposition}
\newtheorem{corollary}[theorem]{Corollary}
\newtheorem{conjecture}[theorem]{Conjecture}
\newtheorem*{mainthm}{Main Theorem}
\newtheorem*{maincor}{Corollary}
\theoremstyle{definition}
\newtheorem{definition}[theorem]{Definition}
\newtheorem{remark}[theorem]{Remark}
\newtheorem{example}[theorem]{Example}

\newcommand{\Z}{\mathbb{Z}}
\newcommand{\Q}{\mathbb{Q}}
\newcommand{\C}{\mathbb{C}}
\newcommand{\R}{\mathcal{R}}
\newcommand{\Aut}{\operatorname{Aut}}
\newcommand{\Out}{\operatorname{Out}}
\newcommand{\Inn}{\operatorname{Inn}}
\newcommand{\End}{\operatorname{End}}
\newcommand{\Sym}{\operatorname{Sym}}
\newcommand{\Hom}{\operatorname{Hom}}
\newcommand{\coker}{\operatorname{coker}}
\newcommand{\Rinf}{R_\infty}
\newcommand{\GAP}[1]{\textbf{[GAP: #1]}}

\title{Property $R_\infty$ for the Artin groups of types $E_6,E_7,E_8,F_4,H_3,H_4$}
\author{[Author names]}
\date{Draft, November 2026}

\begin{document}

\begin{abstract}
A group has property $\Rinf$ if every automorphism has infinitely many twisted conjugacy classes.
We prove that the Artin groups of spherical types $E_6$, $E_7$, $E_8$, $F_4$, $H_3$ and $H_4$ have
property $\Rinf$. Combined with results of Fel'shtyn--Gon\c{c}alves, Calvez--Soroko and Soroko--Vaskou,
this shows that every irreducible Artin group of spherical type and rank at least two has property $\Rinf$.
The proof does not use a description of the automorphism group of the Artin group, which is unknown for
these types. It uses three facts: the pure Artin group $P$ is characteristic (Cohen--Paris);
$H_1(P;\Z)$ is the permutation module $\Z[\R]$ on the reflections; and the induced automorphism of $W/Z(W)$ is
realised by a permutation of $\R$. Schur's lemma, applied to a rational, absolutely irreducible constituent
of multiplicity one of $\Q[\R]$, then gives eigenvalue $1$ on $H_1(P;\Q)$ after twisting by a suitable
inner automorphism. The representation-theoretic input is certified by exact computer calculations, checked
with two independent programs. As by-products we give an independent proof for type $D_5$ and a uniform
proof that the inversion automorphism of any Artin group that is not free abelian has infinite Reidemeister number.
\end{abstract}

\maketitle

\noindent\fbox{\parbox{\dimexpr\textwidth-2\fboxsep-2\fboxrule}{\small
\textbf{AI-assistance statement (placeholder; to be completed by the human authors).}
The mathematical argument, the computer programs and this manuscript were produced with substantial
assistance from an AI system (Anthropic's Claude). They were checked by AI-based review only, and have
\emph{not yet} been checked by a human expert. The authors must replace this paragraph with an accurate
statement of contributions and verification before any submission or posting.}}

\section{Introduction}

\subsection{Twisted conjugacy and property $\Rinf$}
Let $G$ be a group and $\varphi\in\Aut(G)$. Two elements $x,y\in G$ are \emph{$\varphi$-twisted conjugate} if
$y=zx\varphi(z)^{-1}$ for some $z\in G$. The number $R(\varphi)\in\{1,2,\dots\}\cup\{\infty\}$ of
equivalence classes is the \emph{Reidemeister number} of $\varphi$. The group $G$ has \emph{property $\Rinf$}
if $R(\varphi)=\infty$ for every $\varphi\in\Aut(G)$. See \cite{FT15,FN16} for surveys.

For Artin groups the question goes back to Fel'shtyn and Gon\c{c}alves, who proved $\Rinf$ for the braid
groups \cite{FG10}. Juh\'asz treated certain large-type Artin groups \cite{Juh11}. Dekimpe, Gon\c{c}alves and
Ocampo treated the pure braid groups \cite{DGO21}. Calvez and Soroko gave a uniform proof for the spherical
types $A_n,B_n,D_4,I_2(m)$, their pure subgroups, and the affine types $\tilde A_n,\tilde C_n$ \cite{CS22}.
Soroko and Vaskou proved it for spherical type $D_n$, $n\ge 6$ (and for large hyperbolic-type groups)
\cite{SV24}. Soroko proved $\Rinf$ for many non-spherical Artin and Coxeter groups whose outer
automorphisms are length-preserving \cite{Sor26}. The question ``Which Artin groups have property $\Rinf$?''
is Problem~8.6 of the AIM problem list on the geometry and topology of Artin groups \cite{AIMPL}. In each of
the spherical cases above, the proof uses an explicit description of $\Aut(A)$ or $\Aut(A/Z(A))$. For example,
\cite{CS22} embeds $\Aut(A/Z(A))$ in a mapping class group.

\subsection{Main result}
For a Coxeter graph $\Gamma$ let $A=A(\Gamma)$ be the Artin group and $W=W(\Gamma)$ the Coxeter group.

\begin{mainthm}
The Artin groups of types $E_6$, $E_7$, $E_8$, $F_4$, $H_3$ and $H_4$ have property $\Rinf$.
\end{mainthm}

To our knowledge $\Aut(A)$ has not been determined for any of these six types (cf.~the list in
\cite[\S1]{PS26}). Charney and Crisp \cite[Prop.~9]{CC05} obtained partial information for $F_4$; as recorded in
\cite[\S4]{JMS25}, $\Aut(A(F_4))$ is virtually a direct product of infinite groups. Our proof avoids $\Aut(A)$
altogether.

\begin{maincor}
Every irreducible Artin group of spherical type and rank $\ge 2$ has property $\Rinf$.
\end{maincor}

\begin{proof}[Proof (from the literature and the Main Theorem)]
The irreducible spherical types of rank $\ge2$ are $A_n$ $(n\ge2)$, $B_n$ $(n\ge2)$, $D_n$ $(n\ge4)$,
$E_{6,7,8}$, $F_4$, $H_3$, $H_4$ and $I_2(m)$ $(m\ge5)$. Types $A_n,B_n,D_4,I_2(m)$ are covered by \cite{CS22}
(for $A_n$ also \cite{FG10}), and $D_n$ for $n\ge6$ by \cite[Thm.~1]{SV24}. Property $\Rinf$ for $D_5$ is credited
to \cite{SV24} in \cite[\S1]{Sor26}, using the description of $\Aut(A(D_5))$ in \cite{PS26}. (The arXiv version
v2 of \cite{SV24} that we consulted still states $D_5$ as a conjecture; \cite{PS26} itself does not discuss
$\Rinf$.) An independent proof for $D_5$ is given in Proposition~\ref{prop:D5} below. The remaining types are
covered by the Main Theorem.
\end{proof}

Rank one must be excluded: $A(A_1)\cong\Z$, and inversion on $\Z$ has Reidemeister number $2$.

\subsection{What is new, and relation to general criteria}
The new content is the Main Theorem for the six exceptional types and the method, which needs no knowledge of $\Aut(A)$.
Section~\ref{sec:complements} contains an independent proof for $D_5$; this is not a new result. It also
contains a uniform proof that the inversion automorphism has infinite Reidemeister number. That result is not
new in general: for non-spherical types it follows from \cite[Thm.~1]{Sor26}, and see also \cite{Juh11}.

\smallskip\noindent\emph{Conjectures (R) and (NA).} Fel'shtyn and Troitsky asked whether every finitely generated
residually finite group either has $\Rinf$ or is solvable-by-finite \cite{FT15}. This is Conjecture~(R) in
\cite[\S6]{SV24}; see \cite[Problem~19.28]{Kourovka}. A weaker variant is Conjecture~(NA): every finitely
generated, residually finite, non-amenable group has $\Rinf$. The statement (NA) appears as a theorem in the
preprint \cite{FT12}, which would imply our Main Theorem. However, Soroko and Vaskou report that Troitsky has
informed them that the proof contains a mistake \cite[\S6]{SV24}, so (NA) is currently open. Spherical Artin
groups of rank $\ge2$ are finitely generated, linear (hence residually finite) \cite{CW02,Dig03}, and contain
non-abelian free subgroups. Our results therefore verify Conjectures (R) and (NA) for these groups.

\smallskip\noindent\emph{Acylindrical and quasimorphism criteria.} Several recent criteria deduce $\Rinf$ from a
suitable action of $\Aut(G)$, or of $\langle\Inn(G),\varphi\rangle$, on a hyperbolic space, or from
$\Aut$-invariant quasimorphisms \cite{FFW23,FFIMSW25,SV26}. For an irreducible spherical Artin group $A$, the central quotient $A/Z(A)$
is acylindrically hyperbolic \cite{CW17}. However, applying these criteria to $A$ requires control over
$\Aut(A)$ (or $\Aut(A/Z(A))$), for instance an extension of the Calvez--Wiest action to the automorphism group.
This control is not available for types $E,F,H$; see \cite[Remark~4.8 and footnote~3]{JMS25} for a suggestion in
this direction. The $\Rinf$ result of \cite[Cor.~H]{JMS25} needs a separating vertex of the defining graph. In that graph two
generators are joined whenever $m_{st}<\infty$, so for spherical types it is complete. The result also needs a
trivial centre. Both conditions exclude spherical types.

\smallskip\noindent\emph{Method and precedents.} The idea of forcing eigenvalue $1$ on a characteristic abelian
section by means of the representation theory of a reflection group goes back to Dekimpe--Gon\c{c}alves--Ocampo
\cite{DGO21}. They used $\Gamma_2(P_n)/\Gamma_3(P_n)$ and Stembridge's results on eigenvalues \cite{Ste89}. The
group $A/[P,P]$ is the ``crystallographic quotient'' studied by Gon\c{c}alves--Guaschi--Ocampo \cite{GGO17} and by
Marin \cite{Mar15}. The map $f$ of \S\ref{sec:f} embeds $A/[P,P]$ into $\Z^\R\rtimes W$
(Remark~\ref{rem:cryst}). The action of $A$ on $\Z\times\R$ that underlies $f$ already appears in
Cohen--Paris \cite[Lemma~3.7]{CP03}. We did not find these ingredients combined to prove $\Rinf$ for spherical
Artin groups in the literature. Our literature search was limited (several bibliographic APIs were rate-limited),
so we would welcome references.

\subsection{Outline of the proof}
Let $P=\ker(A\to W)$ be the pure Artin group. Let $\R\subset W$ be the set of reflections, and
$\rho\colon W\to\Sym(\R)$ the conjugation action. Let $\varphi\in\Aut(A)$.
\begin{enumerate}[label=(\arabic*)]
\item $\varphi(P)=P$ (Cohen--Paris \cite{CP03}).
\item $H_1(P;\Z)\cong\Z^\R$ as $A$-modules, where $A$ acts through $W$ by $\rho$ with no signs
(Proposition~\ref{prop:H1}). Hence $\varphi_*\in GL(\Z^\R)$ satisfies
$\varphi_*\rho(w)\varphi_*^{-1}=\rho(\varphi_W(w))$.
\item If a permutation matrix $\pi$ realises the automorphism of $\rho(W)\cong W/Z(W)$ induced by $\varphi$,
then $N=\pi^{-1}\varphi_*$ commutes with $\rho(W)$. By Schur's lemma, $N$ acts on a rational, absolutely
irreducible, multiplicity-one constituent $U\subset\Q^\R$ as a rational scalar. Since $N$ preserves the lattice
$U\cap\Z^\R$, the scalar is $\pm1$.
\item After composing $\varphi$ with a suitable inner automorphism, $\varphi_*$ has eigenvalue $1$ on
$H_1(P;\Q)$ (on $U$, or on $U_1\oplus U_2$ in a swap variant used for $F_4$), so
$R(\varphi|_P)=\infty$. The finite-index inequality then gives $R(\varphi)=\infty$.
\end{enumerate}
The automorphisms of $W/Z(W)$ and the constituents $U$ are determined by finite computations. We certify these
by exact arithmetic with two independent programs, except for one input for $E_8$ taken from the ATLAS
(\S\ref{sec:verify}).

\section{Preliminaries}

\subsection{Reidemeister numbers}
\begin{lemma}\label{lem:reid}
Let $G$ be a group and $\psi\in\Aut(G)$.
\begin{enumerate}[label=(\alph*)]
\item If $K\trianglelefteq G$ and $\psi(K)=K$, let $\bar\psi$ be the induced automorphism of $G/K$. Then
$R(\psi)\ge R(\bar\psi)$.
\item If $G\cong\Z^m$ and $\psi$ is given by $B\in GL_m(\Z)$, then $R(\psi)=|\coker(I-B)|$. In particular
$R(\psi)=\infty$ if $B$ has eigenvalue $1$.
\item For $g\in G$ let $\iota_g(x)=gxg^{-1}$. Then $R(\iota_g\circ\psi)=R(\psi)$.
\item If $H\le G$ has finite index and $\psi(H)=H$, then $R(\psi|_H)\le[G:H]\cdot R(\psi)$. In particular, if
$R(\psi|_H)=\infty$ then $R(\psi)=\infty$.
\end{enumerate}
\end{lemma}

\begin{proof}
(a) The projection $p\colon G\to G/K$ satisfies $p\circ\psi=\bar\psi\circ p$. It therefore maps the class of
$x$ into the class of $p(x)$, and it is surjective.

(b) Here $y$ is $\psi$-conjugate to $x$ iff $y-x\in(I-B)\Z^m$. If $\det(I-B)=0$, the cokernel has positive rank.

(c) We have $y=z x(g\psi(z)^{-1}g^{-1})$ iff $yg=z(xg)\psi(z)^{-1}$. So right multiplication by $g$ is a
bijection from the $\iota_g\psi$-classes to the $\psi$-classes.

(d) Write $G=\bigsqcup_{i=1}^k Hg_i$ with $k=[G:H]$. Let $x\in H$, and suppose $y=zx\psi(z)^{-1}\in H$ for some
$z\in G$. Write $z=hg_i$ with $h\in H$. Then $g_ix\psi(g_i)^{-1}=h^{-1}y\psi(h)\in H$, so $y$ is
$\psi|_H$-conjugate to $x_i:=g_ix\psi(g_i)^{-1}$. Thus the $\psi$-class of $x$ meets $H$ in at most $k$
$\psi|_H$-classes. Every $\psi|_H$-class lies in a $\psi$-class, so $R(\psi|_H)\le kR(\psi)$. Normality of
$H$ is not needed. Moreover $\psi(H)\subseteq H$ already implies $\psi(H)=H$, since $\psi(H)$ has the same
finite index as $H$.
\end{proof}

\subsection{Artin groups, Coxeter groups, reflections}
Let $\Gamma$ be a Coxeter graph with vertex set $S$ and Coxeter matrix $(m_{st})$. The Artin group is
$A=\langle S\mid \Pi(s,t;m_{st})=\Pi(t,s;m_{st})\ (s\neq t,\ m_{st}<\infty)\rangle$, where $\Pi(s,t;m)=sts\cdots$
has $m$ letters. The Coxeter group $W$ is $A/\langle\!\langle s^2:s\in S\rangle\!\rangle$. Write
$\mu\colon A\to W$, $a\mapsto\bar a$, and $P=\ker\mu$. Thus $P$ is the normal closure of $\{s^2:s\in S\}$. Let
\[\R=\{wsw^{-1}:w\in W,\ s\in S\}\]
be the set of reflections, and let $\rho\colon W\to\Sym(\R)$, $\rho(w)(r)=wrw^{-1}$. We use the same letter for
the permutation matrices on $\Z^\R$, $\rho(w)e_r=e_{wrw^{-1}}$.

\begin{lemma}\label{lem:rho}
$\ker\rho=Z(W)$, so $\rho(W)\cong W/Z(W)$. Every automorphism $\alpha$ of $W$ preserves $Z(W)$ and induces an
automorphism $\theta_\alpha$ of $\rho(W)$ with $\theta_\alpha(\rho(w))=\rho(\alpha(w))$.
\end{lemma}
\begin{proof}
$\rho(w)=1$ iff $w$ commutes with every reflection. Since $\R$ generates $W$, this holds iff $w\in Z(W)$. The
centre is characteristic, which gives the second statement.
\end{proof}

For the finite irreducible types considered here, $Z(W)=\{1,w_0\}$ has order $2$ when $w_0=-1$ (types
$E_7,E_8,F_4,H_3,H_4$), and $Z(W)=1$ for $E_6$ and $D_5$ \cite[\S3.19]{Hum90}.

\begin{theorem}[Cohen--Paris {\cite[Thm.~1.4]{CP03}}; Franco--Paris {\cite{FP06}}]\label{thm:char}
If $A$ is an irreducible Artin group of spherical type, then $P=\ker(A\to W)$ is a characteristic subgroup of $A$.
Franco and Paris \cite{FP06} extend this to all spherical types, not necessarily irreducible.
\end{theorem}

The statement of \cite[Thm.~1.4]{CP03} reads: ``Let $A$ be an irreducible Artin group of spherical type. Then
$CA$ is a characteristic subgroup of $A$,'' where $CA=\ker(\mu)$ is the coloured Artin group. Thus every
$\varphi\in\Aut(A)$ satisfies $\varphi(P)=P$. Consequently $\varphi$ induces an automorphism $\varphi_W$ of
$W=A/P$, and by Lemma~\ref{lem:rho} an automorphism $\theta_\varphi:=\theta_{\varphi_W}$ of $\rho(W)$.

\subsection{The homomorphism $f$ and the abelianisation of $P$}\label{sec:f}
Let $W$ act on $\Z^\R=\bigoplus_{r\in\R}\Z e_r$ by $w\cdot e_r=e_{wrw^{-1}}$, and form $\Z^\R\rtimes W$ with
multiplication $(v,w)(v',w')=(v+w\cdot v',ww')$.

\begin{lemma}\label{lem:f}
Let $W$ be any Coxeter group. There is a unique homomorphism $f\colon A\to\Z^\R\rtimes W$ with $f(s)=(e_s,s)$ for
$s\in S$. Its second component is $\mu$. For $a\in A$ and a word $a=s_1^{\epsilon_1}\cdots s_\ell^{\epsilon_\ell}$,
\begin{equation}\label{eq:f}
f(a)=\Bigl(\textstyle\sum_{k=1}^{\ell}\epsilon_k\,e_{r_k},\ \bar a\Bigr),\qquad r_k=\bar s_1\cdots\bar s_{k-1}\,s_k\,(\bar s_1\cdots \bar s_{k-1})^{-1}.
\end{equation}
\end{lemma}

\begin{proof}
Note $(e_s,s)^{-1}=(-s\cdot e_s,s)=(-e_s,s)$, since $ss s^{-1}=s$. This gives \eqref{eq:f} by induction on
$\ell$; the sign does not affect $r_k$ because $s^{-1}$ and $s$ have the same image in $W$. We must check the
braid relations. Fix $s\ne t$ with $m=m_{st}<\infty$. Let $u=\Pi(s,t;m)$ and $u'=\Pi(t,s;m)$. The $W$-components agree because the braid relation holds in $W$. In $W$ the element $st$ has order exactly
$m$, and the subgroup $\langle s,t\rangle$ is dihedral of order $2m$ \cite[\S5.3]{Hum90}. By \eqref{eq:f}, the
$\Z^\R$-component of $f(u)$ is $\sum_{k=1}^me_{r_k}$, where $r_k$ is the prefix of length $k-1$ of $u$ conjugating
the $k$-th letter. A direct computation, using $s(ts)^j=(st)^js$, gives $r_k=(st)^{k-1}s$ for both parities of
$k$. Similarly, for $u'=tst\cdots$ the reflections are $r'_k=(ts)^{k-1}t=(ts)^ks=(st)^{m-k}s$. Since $st$ has order
$m$, both $\{r_k\}$ and $\{r'_k\}$ equal $\{(st)^js:0\le j<m\}$, a set of $m$ distinct elements of $\R$. So the
$\Z^\R$-components agree and $f(u)=f(u')$.
\end{proof}

\begin{remark}
Lemma~\ref{lem:f} is equivalent to \cite[Lemma~3.7(1)]{CP03}. There $A$ acts on $\Z\times\R$ by
$U_s(k,t)=(k+\delta_{t,s},sts)$, and this action commutes with the translations $L_k$. Such an action is the same
as a homomorphism to $\Z^\R\rtimes\Sym(\R)$; we include the short proof for completeness. The computation is also
checked by machine (Appendix~\ref{app:comp}).
\end{remark}

\begin{proposition}\label{prop:H1}
Let $A$ be of spherical type. Then $f(P)=2\Z^\R$, and $f|_P$ induces an isomorphism $H_1(P;\Z)\to2\Z^\R$. Under
this isomorphism the conjugation action of $A$ on $H_1(P;\Z)$ becomes $a\cdot x=\rho(\bar a)x$: $\bar a$ permutes
the coordinates, with no signs.
\end{proposition}

\begin{proof}
Since $\mu(P)=1$, $f$ maps $P$ into the abelian group $\Z^\R\times\{1\}$, so $f|_P$ factors through $H_1(P;\Z)$.
We have $f(s^2)=(e_s+s\cdot e_s,1)=(2e_s,1)$. For $a\in A$,
$f(as^2a^{-1})=f(a)(2e_s,1)f(a)^{-1}=(2\bar a\cdot e_s,1)=(2e_{\bar as\bar a^{-1}},1)$. Since $P$ is generated by
the conjugates of the $s^2$, and every reflection is conjugate to a simple one, $f(P)=2\Z^\R$. The same
computation, $f(apa^{-1})=\bar a\cdot f(p)$ for $p\in P$, gives the equivariance.

It remains to show injectivity on $H_1(P;\Z)$. By Brieskorn \cite{Bri71}, $P\cong\pi_1(M)$, where $M=\C^n\setminus\bigcup_{H}H_\C$ is the complement of the
complexified reflection arrangement; one hyperplane $H$ corresponds to each reflection. The first homology of the complement
of a central arrangement of $N$ complex hyperplanes is free abelian of rank $N$, generated by meridians
\cite{Bri73}, \cite[Ch.~5]{OT92}. (Equivalently, the Poincar\'e polynomial of $M$ is $\prod_j(1+m_jt)$, where the
$m_j$ are the exponents, with $\sum_j m_j=|\R|$ \cite{Bri73}, and $H^*(M;\Z)$ is torsion-free \cite[Ch.~5]{OT92}.)
Thus $H_1(P;\Z)\cong\Z^{|\R|}$, and the surjection $H_1(P;\Z)\to2\Z^\R\cong\Z^{|\R|}$ is an isomorphism, because
free abelian groups of finite rank are Hopfian.
\end{proof}

\begin{remark}[Topological interpretation; signs]
Under $P=\pi_1(M)$, the element $as^2a^{-1}$ is a meridian around the hyperplane of the reflection
$\bar as\bar a^{-1}$. The deck group $W$ acts on $M$ by complex-linear maps, which preserve the complex orientation
of normal discs. Hence $W$ permutes the meridian classes without signs, in agreement with
Proposition~\ref{prop:H1}. As a further consistency check, the generator $\Delta^2$ of the pure centre (the full
twist) maps to $2\sum_re_r$; this is verified by machine for all six types.
\end{remark}

\begin{remark}[Crystallographic quotient]\label{rem:cryst}
Since $\ker f\subseteq P$ and $\ker(f|_P)=[P,P]$, $f$ induces an embedding of $A/[P,P]$ into $\Z^\R\rtimes W$ with
image of index $2^{|\R|}$. It contains the lattice $2\Z^\R$ with quotient $W$. This is the crystallographic group of
\cite{GGO17,Mar15}.
\end{remark}

\begin{corollary}\label{cor:twist}
Let $\varphi\in\Aut(A)$ and let $M_\varphi\in GL(\Z^\R)$ be the matrix of the induced automorphism of
$H_1(P;\Z)\cong\Z^\R$. Then
\begin{equation}\label{eq:twist}
M_\varphi\,\rho(w)\,M_\varphi^{-1}=\rho(\varphi_W(w))=\theta_\varphi(\rho(w))\qquad(w\in W),
\end{equation}
and $M_{\iota_a\circ\varphi}=\rho(\bar a)M_\varphi$ for $a\in A$.
\end{corollary}

\begin{proof}
Let $p\in P$ and $a\in A$. Then $\varphi(apa^{-1})=\varphi(a)\varphi(p)\varphi(a)^{-1}$. Passing to $H_1(P)$ and
using Proposition~\ref{prop:H1} gives $M_\varphi\rho(\bar a)x=\rho(\varphi_W(\bar a))M_\varphi x$. Conjugation by
$a$ acts on $H_1(P)$ as $\rho(\bar a)$, which gives the second formula.
\end{proof}

\subsection{Rational, absolutely irreducible constituents of multiplicity one}
\begin{definition}
A $\rho(W)$-invariant subspace $U\subseteq\Q^\R$ is a \emph{good constituent} if $U\otimes\C$ is irreducible
and occurs with multiplicity one in $\C^\R$. Then $U$ is absolutely irreducible, and it is the full isotypic
component of its isomorphism type in $\Q^\R$.
\end{definition}

\begin{lemma}\label{lem:schur}
Let $U$ be a good constituent and let $N\in GL(\Z^\R)$ commute with $\rho(W)$. Then $N(U)=U$ and $N|_U=\pm1$.
\end{lemma}
\begin{proof}
The subspace $N(U)$ is $\rho(W)$-invariant and isomorphic to $U$. By multiplicity one, $N(U)\subseteq U$, hence
$N(U)=U$. Now $N|_U\in\End_{\Q W}(U)=\Q$, since $U$ is absolutely irreducible. So $N|_U=c\in\Q^\times$. The lattice
$\Lambda=U\cap\Z^\R$ has full rank in $U$ because $U$ is defined over $\Q$, and $N^{\pm1}(\Lambda)\subseteq\Lambda$.
Take $v\in\Lambda$ primitive in $\Z^\R$, i.e.\ with coprime coordinates. Then $cv\in\Z^\R$ forces $c\in\Z$, and
likewise $c^{-1}\in\Z$. Hence $c=\pm1$.
\end{proof}

\begin{lemma}[Certificates]\label{lem:cert}
Let $U\subseteq\Q^\R$ be a nonzero $\rho(W)$-invariant subspace.
\begin{enumerate}[label=(\roman*)]
\item The commutant $\End_{\Q W}(\Q^\R)$ is spanned by the orbital matrices $O_\omega$ ($\omega$ a
$\rho(W)$-orbit on $\R\times\R$; $(O_\omega)_{r,r'}=1$ iff $(r,r')\in\omega$). If every $O_\omega$ maps $U$ into
itself and acts on it as a scalar, then $U$ is a good constituent.
\item Let $\chi_U$ be the character of $U$ and $\chi_\R(g)=\#\{r:g(r)=r\}$ the permutation character. If
$\langle\chi_U,\chi_U\rangle=1$ and $\langle\chi_U,\chi_\R\rangle=1$ (inner products over $\rho(W)$), then $U$ is a
good constituent.
\end{enumerate}
\end{lemma}

\begin{proof}
(i) A matrix commutes with all permutation matrices $\rho(w)$ iff its entries are constant on $\rho(W)$-orbits of
$\R\times\R$, so the $O_\omega$ span the commutant. The same holds over $\C$. Write
$\C^\R=\bigoplus_iV_i\otimes\C^{m_i}$ with $V_i$ irreducible and pairwise non-isomorphic. Then
$\End_{\C W}(\C^\R)=\prod_iM_{m_i}(\C)$, acting on the second factors. By hypothesis this algebra acts on
$U_\C=U\otimes\C$ by scalars. The central idempotents act on $U_\C$ by $0$ or $1$, and they sum to $1$. Hence
$U_\C\subseteq V_i\otimes\C^{m_i}$ for a single $i$, and $U_\C=V_i\otimes L$ for a subspace $0\ne L\subseteq\C^{m_i}$.
The algebra $M_{m_i}(\C)$ preserves $L$ and acts on it by scalars, which forces $m_i=1$ and $L=\C$. So $U_\C=V_i$ is
irreducible of multiplicity one.

(ii) The characters of $\rho(W)$ on rational representations are real. So $\langle\chi_U,\chi_U\rangle=1$ means
that $U_\C$ is irreducible, and $\langle\chi_U,\chi_\R\rangle$ is the multiplicity of $U_\C$ in $\C^\R$.
\end{proof}

\section{The criterion}

\begin{theorem}\label{thm:crit}
Let $A$ be an irreducible Artin group of spherical type with Coxeter group $W$. Suppose that for every automorphism
$\alpha$ of $W$ there exists a permutation $\pi$ of $\R$ such that
\begin{equation}\label{eq:pi}
\pi\,\rho(w)\,\pi^{-1}=\theta_\alpha(\rho(w))\quad\text{for all }w\in W,
\end{equation}
and at least one of the following holds.
\begin{enumerate}[label=(\roman*)]
\item\emph{(stable case)} There is a good constituent $U$ such that for each $\varepsilon\in\{1,-1\}$ there exist
$h\in\rho(W)$ and $0\ne v\in U$ with $h\pi v=\varepsilon v$.
\item\emph{(swap case)} There are good constituents $U_1\ne U_2$ with $\pi(U_1)=U_2$ such that for each
$\varepsilon\in\{1,-1\}$ there exist $h\in\rho(W)$ and $0\ne v\in U_1$ with $(h\pi)^2v=\varepsilon v$.
\end{enumerate}
Then $A$ has property $\Rinf$.
\end{theorem}

\begin{proof}
Let $\varphi\in\Aut(A)$. By Theorem~\ref{thm:char}, $\varphi(P)=P$. Apply the hypothesis to $\alpha=\varphi_W$ and
let $\pi$ be the resulting permutation. By \eqref{eq:twist} and \eqref{eq:pi}, $N:=\pi^{-1}M_\varphi$ commutes
with $\rho(W)$. It lies in $GL(\Z^\R)$, since $\pi$ is a permutation matrix.

\emph{Case (i).} By Lemma~\ref{lem:schur}, $N|_U=c\in\{\pm1\}$. Choose $h\in\rho(W)$ and $0\ne v\in U$ with
$h\pi v=cv$, and choose $a\in A$ with $\rho(\bar a)=h$. Put $\chi=\iota_a\circ\varphi$. By
Corollary~\ref{cor:twist},
\[M_\chi v=\rho(\bar a)M_\varphi v=h\pi Nv=c\,h\pi v=c^2v=v.\]
Note that $h\pi$ need not preserve $U$; this is not needed.

\emph{Case (ii).} By Lemma~\ref{lem:schur}, $N|_{U_i}=c_i\in\{\pm1\}$ for $i=1,2$. Choose $h,v$ with
$(h\pi)^2v=c_1c_2v$, $0\neq v\in U_1$, and $a$ with $\rho(\bar a)=h$. Put $\chi=\iota_a\circ\varphi$ and $T=h\pi$,
so $M_\chi=TN$. Since $h$ preserves each isotypic component and $\pi U_1=U_2$, we have $TU_1=U_2$. For $v\in U_1$,
\[M_\chi^2v=TNTNv=c_1\,TNTv=c_1c_2\,T^2v=v .\]
Here $Tv\in U_2$, so $N$ acts on it as $c_2$. Put $u=v+M_\chi v$. Then $M_\chi u=M_\chi v+v=u$. Moreover
$u\ne0$: $v\in U_1\setminus\{0\}$, $M_\chi v=c_1Tv\in U_2$, and $U_1\cap U_2=0$, since distinct good constituents
are distinct isotypic components.

In both cases $M_\chi$ has eigenvalue $1$ on $\Q^\R$, hence a nonzero fixed vector in $\Z^\R$ after clearing
denominators. By Lemma~\ref{lem:reid}(b) and Proposition~\ref{prop:H1}, the automorphism of
$H_1(P;\Z)=P/[P,P]$ induced by $\chi|_P$ has infinite Reidemeister number. The subgroup $[P,P]$ is characteristic
in $P$, so Lemma~\ref{lem:reid}(a) gives $R(\chi|_P)=\infty$. By Lemma~\ref{lem:reid}(d), with
$[A:P]=|W|<\infty$, we get $R(\chi)=\infty$. By Lemma~\ref{lem:reid}(c), $R(\varphi)=R(\chi)=\infty$.
\end{proof}

\begin{remark}[Reduction to outer classes]\label{rem:reduce}
The hypotheses depend only on the class of $\theta_\alpha$ in $\Out(\rho(W))$. If $\pi$ satisfies \eqref{eq:pi}
for $\theta$ and $g\in\rho(W)$, then $g\pi$ satisfies it for $\iota_g\circ\theta$. Moreover
$\{hg\pi:h\in\rho(W)\}=\{h\pi\}$, so (i) and (ii) for $g\pi$ are equivalent to (i) and (ii) for $\pi$. Hence it
suffices to treat one representative of each class of $\Out(\rho(W))$ that is induced by $\Aut(W)$. For the trivial
class take $\pi=1$. Then (i) holds as soon as $U$ is a good constituent on which some reflection acts
nontrivially: $h=1$ gives eigenvalue $+1$, and an involution acting nontrivially has eigenvalue $-1$.
\end{remark}

\begin{corollary}\label{cor:inner}
Suppose every automorphism of $\rho(W)$ is inner, and $\Q^\R$ has a good constituent on which some simple
reflection acts nontrivially. Then $A$ has property $\Rinf$.
\end{corollary}

\begin{remark}[How permutations $\pi$ arise]\label{rem:pi-refl}
Suppose an automorphism $\theta$ of $\rho(W)$ maps the set $\rho(\R)$ of reflections onto itself, and
$r\mapsto\rho(r)$ is injective on $\R$ (true in rank $\ge3$, and checked by machine in all cases). Then
$\pi(r):=\rho^{-1}\theta(\rho(r))$ satisfies \eqref{eq:pi}. Indeed,
\[\pi\rho(w)\pi^{-1}(r)=\rho^{-1}\theta\bigl(\rho(w)\,\theta^{-1}(\rho(r))\,\rho(w)^{-1}\bigr)=\rho^{-1}\bigl(\theta(\rho(w))\,\rho(r)\,\theta(\rho(w))^{-1}\bigr)=\theta(\rho(w))(r).\]
\end{remark}

\section{Verification for $E_6,E_7,E_8,F_4,H_3,H_4$}\label{sec:verify}

\subsection{Automorphisms of $W/Z(W)$}
\begin{proposition}\label{prop:out}
Let $G=\rho(W)\cong W/Z(W)$.
\begin{enumerate}[label=(\alph*)]
\item For $E_6$, $E_7$ and $D_5$, $\Out(G)=1$.
\item For $H_3$, $H_4$ and $F_4$, $|\Out(G)|=2$, and the nontrivial class is represented by conjugation with the
explicit permutation $\pi_0$ of $\R$ described below.
\item For $E_8$, $\Out(G)=1$. This relies on the ATLAS (see the proof).
\end{enumerate}
\end{proposition}

\begin{proof}
(a), (b): by computer, in two independent ways (Appendix~\ref{app:comp}).
\begin{enumerate}[label=(\arabic*)]
\item \texttt{fastaut.py}/\texttt{pgrp.py} count all automorphisms by brute force over images of the Coxeter
generators. The candidate images are involutions in conjugacy classes of the right size, subject to the Coxeter
relations. The program checks that the map extends to a homomorphism, verifying the relations on a spanning tree
of the Cayley graph, and that it is bijective. Results: $|\Aut G|=|G|$ for $E_6$ ($51840$) and $D_5$ ($1920$);
$|\Aut G|=2|G|$ for $H_3$ ($120$), $F_4$ ($1152$) and $H_4$ ($14400$).
\item \texttt{indep/out\_check.py} is fresh code using exact root systems. It enumerates $G$, its involutions and
their conjugacy classes. It checks that no involution class other than a reflection class has the size of a
reflection class, so every automorphism maps $\rho(\R)$ onto itself. Let $(t_1,\dots,t_n)$ be a generating tuple of reflections with
$\mathrm{ord}(t_it_j)=\mathrm{ord}(\rho(s_i)\rho(s_j))$, which divides $m_{ij}$. Such a tuple defines a surjection
$W\to G$, $s_i\mapsto t_i$. Its kernel has order $|Z(W)|\le 2$, so it is
central and equals $Z(W)$. Such tuples therefore correspond bijectively to automorphisms. The program counts their
$\Inn(G)$-orbits by a stabiliser-chain search. The results agree with (1), and also give $\Out(G)=1$ for $E_7$
($|G|=1451520$; involution classes of sizes $63,315,945,3780$).
\end{enumerate}

\emph{The permutations $\pi_0$.} Roots are normalised as in \S\ref{app:roots}.
\begin{itemize}
\item $H_3$, $H_4$: roots have coordinates in $\Z[\phi]$, $\phi=(1+\sqrt5)/2$. Let $\sigma$ be the Galois
automorphism $\sqrt5\mapsto-\sqrt5$, applied coordinatewise. Then $T=\tau\circ\sigma$ maps the root system $\Phi$
onto itself, where $\tau$ swaps the last two coordinates. $T$ is $\sigma$-semilinear and satisfies
$(Tx,Ty)=\sigma((x,y))$, hence $Ts_\alpha T^{-1}=s_{T\alpha}$. The induced permutation $\pi_0$ of $\R$ normalises
$\rho(W)$. The induced automorphism is outer; this is checked directly by machine, by verifying that no element of
$\rho(W)$ induces the same automorphism. Moreover $\pi_0^2=1$, since $\sigma$ and $\tau$ commute and are involutions.
\item $F_4$: $L(x)=(x_1+x_2,x_1-x_2,x_3+x_4,x_3-x_4)$ is $\sqrt2$ times an orthogonal map. It sends short roots to
long roots and long roots to twice short roots, so $Ls_\alpha L^{-1}=s_{L\alpha}$. The induced permutation $\pi_0$
of $\R$ normalises $\rho(W)$ and exchanges the two reflection classes, so it is outer. Since $L^2=2I$, $\pi_0^2=1$.
\end{itemize}

(c) We have $Z(W(E_8))=\{\pm1\}$ and $W(E_8)/\{\pm1\}\cong O_8^+(2){:}2$, the full orthogonal group of the
quadratic form $q(x)=\tfrac12(x,x)\bmod2$ on $L/2L$, $L$ the $E_8$ lattice \cite[p.~85]{ATLAS}; see also
\cite[\S2.12]{Hum90}. The group $S=O_8^+(2)$ (simple, often written $\Omega_8^+(2)$ or $D_4(2)$) has
$\Out(S)\cong S_3$ \cite[p.~85]{ATLAS}. Let $G=S.2\le\Aut(S)$ with $C_G(S)=1$. Then $S$ is the unique minimal
normal subgroup of $G$, every automorphism of $G$ restricts to an automorphism of $S$, and restriction gives
$\Aut(G)\cong N_{\Aut(S)}(G)$. (If $\alpha|_S=1$ then $\alpha(g)g^{-1}\in C_G(S)=1$.) Hence
$\Out(G)\cong N_{S_3}(C_2)/C_2=1$. Independently, \texttt{indep/out\_check\_stab.py} verifies by Schreier--Sims
stabiliser chains that $|G|=348364800$ and that every automorphism of $G$ preserving $\rho(\R)$ is inner. It does
not verify that every automorphism preserves $\rho(\R)$; for $E_8$ this step rests on the ATLAS.
\GAP{The ATLAS page references (p.~85 for $O_8^+(2)$; pp.~26, 46 for $U_4(2)$, $S_6(2)$) were given from memory
and have not been checked against a copy of the ATLAS. For $E_6$ and $E_7$ the computation makes the citation
unnecessary.}
\end{proof}

\begin{remark}
For $E_6$ and $E_7$ the results agree with the classical identifications
$W(E_6)\cong U_4(2){:}2\cong O_6^-(2)$ and $W(E_7)\cong 2\times S_6(2)$, with $\Out(U_4(2){:}2)=\Out(S_6(2))=1$
\cite{ATLAS}. The involution class sizes of $W(E_7)/Z$ found by the program ($63,315,945,3780$) are those of
$S_6(2)$.
\end{remark}

In each case $\Out(G)$ is computed for the whole group $\Aut(\rho(W))$, which contains all automorphisms
$\theta_\alpha$ induced from $\Aut(W)$. By Remarks~\ref{rem:reduce} and~\ref{rem:pi-refl}, the
existence of $\pi$ in \eqref{eq:pi} is settled for every $\alpha\in\Aut(W)$: take $\pi=g$ or $\pi=g\pi_0$ with
$g\in\rho(W)$.

\subsection{Good constituents and witnesses}
\begin{proposition}\label{prop:data}
Table~\ref{tab:data} lists, for each type, good constituents of $\Q^\R$ and the decomposition of $\C^\R$. On each
listed constituent every simple reflection acts nontrivially (for $F_4$, $U_{\mathrm{long}}$ is moved by the long
simple reflections and $U_{\mathrm{short}}$ by the short ones). Moreover:
\begin{itemize}
\item $H_3$, $H_4$: $\pi_0$ preserves each listed $U$, and $\pi_0|_U$ has both eigenvalues $+1$ and $-1$.
So (i) of Theorem~\ref{thm:crit} holds for the outer class with $h=1$ for both signs.
\item $F_4$: $\pi_0(U_{\mathrm{short}})=U_{\mathrm{long}}$ and vice versa. Also $(1\cdot\pi_0)^2=1$ has
eigenvalue $+1$ on $U_{\mathrm{short}}$, and $(s_1\pi_0)^2$ has eigenvalue $-1$ on $U_{\mathrm{short}}$
($s_1$ the long simple reflection $e_2-e_3$). So (ii) holds.
\end{itemize}
\end{proposition}

\begin{table}[ht]
\centering\small
\begin{tabular}{lrrrrll}
\toprule
type & $|W|$ & $|Z(W)|$ & $|\R|$ & $|\rho(W)|$ & $|\Out\rho(W)|$ & $\C^\R$ (good constituents in bold)\\
\midrule
$H_3$ & 120 & 2 & 15 & 60 & 2 & $1+\mathbf 4+2\cdot 5$\\
$H_4$ & 14400 & 2 & 60 & 7200 & 2 & $1+9+9'+\mathbf{16}+\mathbf{25}$\\
$F_4$ & 1152 & 2 & 24 & 576 & 2 & $2\cdot1+\mathbf{2}_{\rm short}+\mathbf 2_{\rm long}+2\cdot 9$\\
$E_6$ & 51840 & 1 & 36 & 51840 & 1 & $1+\mathbf{15}+\mathbf{20}$\\
$E_7$ & 2903040 & 2 & 63 & 1451520 & 1 & $1+\mathbf{27}+\mathbf{35}$\\
$E_8$ & 696729600 & 2 & 120 & 348364800 & 1 (ATLAS) & $1+\mathbf{35}+\mathbf{84}$\\
\midrule
$D_5$ & 1920 & 1 & 20 & 1920 & 1 & $1+\mathbf4+\mathbf5+\mathbf{10}$\\
\bottomrule
\end{tabular}
\caption{Data for the verification. Good constituents were certified exactly by \texttt{certify.py} (commutant
criterion, Lemma~\ref{lem:cert}(i)) and, independently, by \texttt{indep/indep\_check.py} (character criterion,
Lemma~\ref{lem:cert}(ii)) for all types except $E_8$. For $E_8$ the second check uses the commutant criterion with
fresh code: the action of $\rho(W(E_8))$ on $\R$ has rank $3$. In the $F_4$ row, the $2\cdot 1$ and $2\cdot9$ are
read off from character norms ($\langle\chi,\chi\rangle=4$).}
\label{tab:data}
\end{table}

\begin{proof}[Proof of Proposition~\ref{prop:data}]
Exact computer calculation (Appendix~\ref{app:comp}). Each $U$ is the kernel, over $\Q$, of an integer matrix
lying in the commutant, so $U$ is defined over $\Q$. It is certified as a good constituent by
Lemma~\ref{lem:cert}(i) (\texttt{certify.py}) and by Lemma~\ref{lem:cert}(ii) (\texttt{indep\_check.py}). The
eigenvalue statements are certified by computing exact rational kernels: $\dim\ker((h\pi_0)^{(2)}-\varepsilon)\cap
U>0$. This is done by \texttt{certify\_outer.py} for the abstractly computed outer automorphism, and by
\texttt{indep\_check.py} for the explicit $\pi_0$. The two programs find, respectively, the dimensions
$(3,1)$ for $H_3$ and $(10,6)$, $(15,10)$ for $H_4$ (for $\varepsilon=+1,-1$). For $F_4$ the witness eigenspaces
have dimension $1$, or $2$ when $h=1$ and $\varepsilon=+1$, since $\pi_0^2=1$. For $H_3$ and $H_4$ the
\texttt{certify\_outer.py} witnesses were found for some $h$. For the explicit $\pi_0$, $h=1$ already works.
\end{proof}

\begin{proof}[Proof of the Main Theorem]
$E_6$, $E_7$, $E_8$: by Proposition~\ref{prop:out}, $\Out(\rho(W))=1$. Table~\ref{tab:data} gives a good
constituent moved by a simple reflection. Apply Corollary~\ref{cor:inner}.

$H_3$, $H_4$, $F_4$: $\Out(\rho(W))$ has order $2$. The trivial class is handled as in
Remark~\ref{rem:reduce}, with any listed $U$. The nontrivial class is handled by $\pi_0$ and
Proposition~\ref{prop:data}: case (i) for $H_3$, $H_4$ and case (ii) for $F_4$. By Remark~\ref{rem:reduce} this
covers every $\alpha\in\Aut(W)$. Apply Theorem~\ref{thm:crit}.
\end{proof}

\section{Complements}\label{sec:complements}

\begin{proposition}[Type $D_5$; independent proof of a known result]\label{prop:D5}
$A(D_5)$ has property $\Rinf$.
\end{proposition}
\begin{proof}
$\Out(\rho(W(D_5)))=1$ by computer (both programs). This is consistent with the fact that the diagram involution
of $D_5$ is induced by the sign change of $e_5$, and conjugation by it agrees with conjugation by $-\mathrm{diag}(1,1,1,1,-1)=\mathrm{diag}(-1,-1,-1,-1,1)\in W(D_5)$. The constituents of dimensions $4,5,10$ are
good and are moved by every simple reflection (Table~\ref{tab:data}). Apply Corollary~\ref{cor:inner}.
\end{proof}

\begin{remark}[Centre-fixing automorphisms are elementary]\label{rem:centre}
Let $A$ be irreducible spherical. Then $Z(A)=\langle z\rangle\cong\Z$, with $z=\Delta$ or $\Delta^2$
\cite{BS72,Del72}. Hence $\varphi(z)=z^{\pm1}$ for every $\varphi\in\Aut(A)$. The image of $\Delta$ in
$H_1(A;\Z)$ is nonzero: its coordinate sum is $\ell(w_0)=|\R|$. So if $\varphi(z)=z$, the induced automorphism of
$H_1(A;\Z)$ has eigenvalue $1$, and $R(\varphi)=\infty$ by Lemma~\ref{lem:reid}(a,b). Thus only centre-inverting
automorphisms need Theorem~\ref{thm:crit}. For $E_{6,7,8}$ and $H_{3,4}$, where $H_1(A)=\Z$, these are exactly
the automorphisms acting by $-1$ on $H_1(A)$. The same conclusion also follows from the theorem of
Fel'shtyn--Luchnikov--Troitsky \cite{FLT15}, as quoted in \cite[\S1]{Tro22}: if $G$ is finitely generated and
residually finite and $R(\varphi)<\infty$, then $\varphi$ has finitely many fixed points, building on Jabara
\cite{Jab08}. \GAP{We have only seen this statement as quoted in \cite{Tro22}, not in \cite{FLT15} itself. The
remark is not used in the proof.}
\end{remark}

\begin{proposition}[Inversion; not new in general]\label{prop:inv}
Let $A$ be an Artin group that is not free abelian, and let $\iota\in\Aut(A)$ be given by $\iota(s)=s^{-1}$ for
$s\in S$. Then $R(\iota)=\infty$.
\end{proposition}
\begin{proof}
$\iota$ is an automorphism: it maps the relation $\Pi(s,t;m)=\Pi(t,s;m)$ to the inverse of the relation obtained by
reversing both words, which is again a defining relation. Let $f$ be as in Lemma~\ref{lem:f}, valid for any
Coxeter group, and let $\tau(v,w)=(-v,w)$. This is an automorphism of $\Z^\R\rtimes W$, and $f\circ\iota=\tau\circ f$
on generators, since $f(s^{-1})=(-e_s,s)$. So $f$ maps $\iota$-classes into $\tau$-classes. For $g=(u,x)$ and
$y=(v,w)$ one computes $gy\tau(g)^{-1}=(u+xv+xwx^{-1}u,\ xwx^{-1})$. Fix $s,t\in S$ with $m_{st}\ge3$; they exist
because $A$ is not free abelian. Then $sts\ne t$. On elements with $W$-part $s$, twisted conjugation acts by
$v\mapsto(1+s)u+xv$ with $x\in C_W(s)$. Define $D_v(r)=v_r-v_{srs}$. The term $(1+s)u$ does not change $D$, and
$x\in C_W(s)$ transforms $D_v$ by the permutation $r\mapsto x^{-1}rx$. Hence
$\max_r|D_v(r)|$ is an invariant of the $\tau$-class. For $a_k=t^{2k}s$ we have $f(a_k)=(2ke_t+e_s,s)$ and
$\max_r|D(r)|=2k$. So the $a_k$ ($k\ge1$) lie in pairwise distinct $\iota$-classes.
\end{proof}

\begin{remark}[Scope]
The method needs a characteristic finite-index subgroup with a computable abelianisation; the pure subgroup and
$\Z^\R$ play this role. It applies verbatim to the other irreducible spherical types once $\Out(W/Z(W))$ and a good
constituent are known. We have not attempted to extend it to the pure subgroups $P$ themselves, to the central
quotients $A/Z(A)$, or to reducible types. For reducible types the statement can fail, since $A(A_1\times A_1)=\Z^2$.
\end{remark}

\appendix
\section{Computational methods and certificates}\label{app:comp}

All code is in the accompanying package (\texttt{code/}); exact commands and outputs are in
\texttt{code/README.md}. Arithmetic is exact (Python \texttt{fractions}, SymPy \texttt{DomainMatrix} over $\Q$)
wherever a certificate is claimed. Floating point is used only to \emph{locate} candidate eigenvalues or
witnesses, which are then re-checked exactly.

\subsection{Root systems}\label{app:roots}
Original code (\texttt{roots.py}) builds roots in floating point and identifies them after rounding. The
independent code (\texttt{indep/rootsys.py}) builds them exactly. Types $E_{6,7,8}$ use the Bourbaki simple roots
$\alpha_1=\frac12(1,-1,\dots,-1,1)$, $\alpha_2=e_1+e_2$, $\alpha_k=e_{k-2}-e_{k-3}$; $F_4$ uses $e_2-e_3,e_3-e_4,e_4,
\frac12(e_1-e_2-e_3-e_4)$; $D_5$ uses $e_i-e_{i+1}$ and $e_4+e_5$. For $H_3$ the roots are the $30$ vectors
$\pm2e_i$ and the even coordinate permutations of $(\pm\phi,\pm1,\pm\phi^{-1})$; for $H_4$ the $120$ vectors
$\pm2e_i$, $(\pm1,\pm1,\pm1,\pm1)$ and the even permutations of $(\pm\phi,\pm1,\pm\phi^{-1},0)$. These have
coordinates in $\Q(\sqrt5)$, and the simple roots are found by search from the Coxeter Gram matrix. In each case the
program checks the root count and recovers the Coxeter matrix.

\subsection{Certificates}
\begin{enumerate}[label=(\arabic*)]
\item \texttt{certify.py} (original). It computes the orbital matrices $O_\omega$ and a symmetric integer
combination $X=\sum c_\omega(O_\omega+O_\omega^{T})$, and for each candidate integer eigenvalue $\lambda$ the exact
kernel $U=\ker_\Q(X-\lambda)$. It checks exactly that every $O_\omega$ acts on $U$ as a scalar
(Lemma~\ref{lem:cert}(i)), and which simple reflections act nontrivially. Output: Table~\ref{tab:data}, good
constituents of dimensions $4$ $(H_3)$; $2,2$ $(F_4)$; $16,25$ $(H_4)$; $15,20$ $(E_6)$; $27,35$ $(E_7)$; $35,84$
$(E_8)$; $4,5,10$ $(D_5)$.
\item \texttt{certify\_outer.py} (original). It computes an outer automorphism of $\rho(W)$ by brute force and the
permutation $\pi$ of Remark~\ref{rem:pi-refl}, and checks \eqref{eq:pi} on generators. It then certifies the
witnesses by the exact dimension of $\ker(X-\lambda)\cap\ker(Q-\varepsilon)$, with $Q=\rho(h)\pi$ or $(\rho(h)\pi)^2$.
\item \texttt{indep/indep\_check.py} (fresh). It enumerates $\rho(W)$ and checks its order and trivial centre. It
checks the braid-relation computation of Lemma~\ref{lem:f}. It computes joint eigenspaces of the reflection class
sums $\sum_{r\in K}\rho(r)$, which are central elements of the group algebra, and certifies them by exact
character inner products (Lemma~\ref{lem:cert}(ii)). The characters are computed from the exact rational
orthogonal projection. It then checks the explicit $\pi_0$ (normalises, outer, image of constituents, exact
witnesses). For $E_8$ it reruns the commutant criterion with fresh code.
\item \texttt{indep/out\_check.py}, \texttt{indep/out\_check\_stab.py} (fresh): Proposition~\ref{prop:out}.
\item \texttt{check.py} (original). Numerical sanity checks: braid relations for $f$, $f(\Delta^2)=2\sum e_r$, and
random tests of the invariant in Proposition~\ref{prop:inv}.
\end{enumerate}

\subsection{Runtimes}
On a shared, heavily loaded machine, single-threaded: \texttt{certify.py} (all seven types) $\approx 2$\,s;
\texttt{certify\_outer.py H3 F4 H4} $\approx 9$\,s; \texttt{check.py} $\approx12$\,s; \texttt{fastaut.py E6}
$\approx 2.4$\,min (an earlier run under heavier load took $24$\,min); \texttt{indep\_check.py} (H3, F4, H4, E6, E7)
$\approx 9$\,s; \texttt{out\_check.py} (incl.\ $E_7$) $\approx 8$\,s; \texttt{out\_check\_stab.py E7 E8}
$\approx 2$\,s.

\noindent\GAP{Section and page numbers quoted for \cite{Hum90} (\S\S2.12, 3.19, 5.3), \cite{OT92} (Ch.~5) and \cite{ATLAS} (pp.~26, 46, 85) are given from memory and must be checked against the sources. All arXiv statements quoted in the text (\cite{CP03,FP06,SV24,Sor26,PS26,JMS25,Tro22,DGO21}) were checked against the arXiv versions.}

\begin{thebibliography}{99}
\bibitem{AIMPL} AIM Problem List, \emph{Geometry and topology of Artin groups}, Problem 8.6, aimpl.org/geomartingp/8.
\bibitem{ATLAS} J.\,H. Conway, R.\,T. Curtis, S.\,P. Norton, R.\,A. Parker, R.\,A. Wilson, \emph{ATLAS of Finite Groups}, Oxford Univ.\ Press, 1985.
\bibitem{Bri71} E. Brieskorn, Die Fundamentalgruppe des Raumes der regul\"aren Orbits einer endlichen komplexen Spiegelungsgruppe, \emph{Invent.\ Math.} 12 (1971), 57--61.
\bibitem{Bri73} E. Brieskorn, Sur les groupes de tresses [d'apr\`es V.\,I. Arnol'd], S\'eminaire Bourbaki 1971/72, Exp.~401, Lecture Notes in Math.\ 317, Springer, 1973, 21--44.
\bibitem{BS72} E. Brieskorn, K. Saito, Artin-Gruppen und Coxeter-Gruppen, \emph{Invent.\ Math.} 17 (1972), 245--271.
\bibitem{CS22} M. Calvez, I. Soroko, Property $R_\infty$ for some spherical and affine Artin--Tits groups, \emph{J.\ Group Theory} 25 (2022), 1045--1054; arXiv:2106.13260.
\bibitem{CW17} M. Calvez, B. Wiest, Acylindrical hyperbolicity and Artin--Tits groups of spherical type, \emph{Geom.\ Dedicata} 191 (2017), 199--215.
\bibitem{CC05} R. Charney, J. Crisp, Automorphism groups of some affine and finite type Artin groups, \emph{Math.\ Res.\ Lett.} 12 (2005), 321--333.
\bibitem{CP03} A.\,M. Cohen, L. Paris, On a theorem of Artin, \emph{J.\ Group Theory} 6 (2003), 421--441; arXiv:math/0204242.
\bibitem{CW02} A.\,M. Cohen, D.\,B. Wales, Linearity of Artin groups of finite type, \emph{Israel J.\ Math.} 131 (2002), 101--123.
\bibitem{DGO21} K. Dekimpe, D.\,L. Gon\c{c}alves, O. Ocampo, The $R_\infty$ property for pure Artin braid groups, \emph{Monatsh.\ Math.} 195 (2021), 15--33; arXiv:2006.08286.
\bibitem{Del72} P. Deligne, Les immeubles des groupes de tresses g\'en\'eralis\'es, \emph{Invent.\ Math.} 17 (1972), 273--302.
\bibitem{Dig03} F. Digne, On the linearity of Artin braid groups, \emph{J.\ Algebra} 268 (2003), 39--57.
\bibitem{FG10} A. Fel'shtyn, D.\,L. Gon\c{c}alves, Twisted conjugacy classes in symplectic groups, mapping class groups and braid groups (with an appendix by F.~Dahmani), \emph{Geom.\ Dedicata} 146 (2010), 211--223.
\bibitem{FLT15} A. Fel'shtyn, N. Luchnikov, E. Troitsky, Twisted inner representations, \emph{Russ.\ J.\ Math.\ Phys.} 22 (2015), 301--306.
\bibitem{FN16} A. Fel'shtyn, T. Nasybullov, The $R_\infty$ and $S_\infty$ properties for linear algebraic groups, \emph{J.\ Group Theory} 19 (2016), 901--921.
\bibitem{FT12} A. Fel'shtyn, E. Troitsky, Twisted conjugacy classes in residually finite groups, preprint, arXiv:1204.3175.
\bibitem{FT15} A. Fel'shtyn, E. Troitsky, Aspects of the property $R_\infty$, \emph{J.\ Group Theory} 18 (2015), 1021--1034.
\bibitem{FFW23} F. Fournier-Facio, R.\,D. Wade, Aut-invariant quasimorphisms on groups, arXiv:2211.00800.
\bibitem{FFIMSW25} F. Fournier-Facio, H. Iveson, A. Martino, W. Sgobbi, P. Wong, Property $R_\infty$ for groups with infinitely many ends, arXiv:2504.12002.
\bibitem{FP06} N. Franco, L. Paris, On a theorem of Artin, II, \emph{J.\ Group Theory} 9 (2006), 731--751; arXiv:math/0411193.
\bibitem{GGO17} D.\,L. Gon\c{c}alves, J. Guaschi, O. Ocampo, Quotients of the Artin braid groups and crystallographic groups, arXiv:1503.04527 (journal version in J.\ Algebra).
\bibitem{Hum90} J.\,E. Humphreys, \emph{Reflection Groups and Coxeter Groups}, Cambridge Univ.\ Press, 1990.
\bibitem{Jab08} E. Jabara, Automorphisms with finite Reidemeister number in residually finite groups, \emph{J.\ Algebra} 320 (2008), 3671--3679.
\bibitem{JMS25} O. Jones, G. Mangioni, G. Sartori, JSJ splittings for all Artin groups, arXiv:2506.17157.
\bibitem{Juh11} A. Juh\'asz, Twisted conjugacy in certain Artin groups, in: \emph{Ischia Group Theory 2010}, World Sci., 2012, 175--195.
\bibitem{Kourovka} E.\,I. Khukhro, V.\,D. Mazurov (eds.), \emph{The Kourovka Notebook: Unsolved Problems in Group Theory}, 21st ed., Novosibirsk, 2026.
\bibitem{Mar15} I. Marin, Crystallographic groups and flat manifolds from complex reflection groups, arXiv:1509.06213.
\bibitem{OT92} P. Orlik, H. Terao, \emph{Arrangements of Hyperplanes}, Grundlehren 300, Springer, 1992.
\bibitem{PS26} L. Paris, I. Soroko, Automorphisms of the Artin group of type $D_5$, arXiv:2606.26082.
\bibitem{Sor26} I. Soroko, Property $R_\infty$ for Artin and Coxeter groups whose outer automorphisms are length-preserving, arXiv:2610.13904.
\bibitem{SV24} I. Soroko, N. Vaskou, Property $R_\infty$ for new classes of Artin groups, arXiv:2409.18123 (v2, 2026).
\bibitem{SV26} I. Soroko, N. Vaskou, Acylindrical hyperbolicity of automorphism groups and property $R_\infty$ for generalized Higman groups, arXiv:2604.27526.
\bibitem{Ste89} J.\,R. Stembridge, On the eigenvalues of representations of reflection groups and wreath products, \emph{Pacific J.\ Math.} 140 (1989), 353--396.
\bibitem{Tro22} E. Troitsky, Twisted conjugacy in residually finite groups of finite Pr\"ufer rank, arXiv:2210.00591.
\end{thebibliography}

\end{document}
